\subsection{Required-Skill Labels and Task Types}
\label{app:data-labels}
Each instance has a task identifier and a nonempty annotated set of required skill identifiers. We use these released associations as the routing targets, preserving all labels of a retained task. TheoremQA concerns theorem application, LogicBench logical reasoning, ToolQA tool-use workflows, MedCalc-Bench medical calculators, CHAMP mathematical concepts, and BigCodeBench software-library use. All retained TheoremQA, LogicBench, ToolQA, and MedCalc-Bench tasks have exactly one annotated required skill. BigCodeBench and CHAMP also test recovery of multiple required skills.

All 1,059 retained BigCodeBench tasks have multiple labels. CHAMP contains 90 single-skill and 78 multi-skill tasks. Across the six collections, the retained data contain 6,932 task--skill annotation occurrences, distinct from the 306 unique skills. The test split contains 1,097 annotation occurrences over 748 tasks: 608 tasks have one required skill and 140 have multiple required skills. The latter comprise all 123 BigCodeBench test tasks and 17 of the 26 CHAMP test tasks. BigCodeBench test tasks require between two and six annotated skills; CHAMP test tasks require between one and four.

Routing recall uses the entire annotated required set as its denominator, so recovering one relevant skill does not necessarily recover all requirements of a multi-skill task. For annotation-based precision and F1, selected skills outside this set count as false positives. This convention defines the evaluation target; it does not imply that an unannotated skill cannot help a downstream agent. We do not replace a multi-skill target with a single representative label.
